\chapter{Harvard Professor Explains The Rules of Writing}

この講義は、実践的な約束から始まる。私たちは、言語の理論そのものへ入るよう求められているのではない。うまく書くにはどうすればよいのか、そして、大規模言語モデルがすでに滑らかで手堅い散文を生み出せる時代にあって、どうすれば書き続けてうまく書けるのかを問われているのである。Pinker は、この問いに、小手先の助言の寄せ集めではなく、因果的主張、対比、そして実演例の連鎖によって答える。この講義には文字どおりの黒板数学はない。したがって以下の数式は、散文の中で述べられた機構を図式的に再構成したものである。論理を鋭く保ちながら transcript に忠実であるためのものだ。

\section{悪文と知識の呪い}

聞き手はただちに、この講義を支配する問題を開く。なぜ、これほど多くの悪文があるのか。Pinker はまず、もっとも気分のよい誤答を払いのける。分かりにくい散文は、おもに戦略的なものだ、と言いたくなる。つまり書き手たちは、薄い考えを jargon の背後に隠すことによって、洗練を装い、仲間内の結束を守り、部外者を締め出しているのだ、と。Pinker は、それが起こりうること自体は否定しない。だが、それが主たる説明ではない、と言い張る。彼は、読むに堪えない prose を生み出しながら、なお語るべきことを十分に持った、あまりにも多くの優秀な人々を知っているからである。

この反転は重要である。なぜなら、それによって講義が正しい軌道に乗るからだ。中心的な失敗は、道徳的な芝居ではない。較正の失敗なのである。

\begin{definition}
The curse of knowledge とは、自分自身が知っていることを知らないとはどういうことかを、再構成することの難しさである。
\end{definition}

Pinker は、この失敗を、書き手が本来なら問うべきだが、しばしば間に合わない問いへ圧縮する。聴衆は何を知っていて、何を知らないのか。これを図式的な機構へ変えれば、
\begin{equation}
\text{専門知} \to \text{読者をモデル化する失敗} \to \text{jargon / acronyms / abstractions} \to \text{読者の混乱}.
\end{equation}
となる。

\begin{remark}
この表示は編集上の略記であって、講義中に可視化されていた記法ではない。その目的は、Pinker の因果連鎖を明示することに尽きる。
\end{remark}

講義はそのあと、例によってこの式を正当化する。卓越した分子生物学者が、まるで相手も分子生物学者であるかのように TED talk を行う。彼は問題から始めない。なぜその研究が重要なのかも説明しない。実験へまっすぐ突入する。会場はほとんどたちまち置き去りにされるのだが、それに気づかない唯一の人物が、話し手本人である。Pinker の要点は、厳しいが精密である。その男は愚かではない。ただし、一局所的で破滅的な仕方においてだけ例外である。彼は、自分が知っていることを知らないとはどんなことかを知らないのだ。

同じ失敗は、抽象が説明を追い越すときにも現れる。Pinker は、``the level of the stimulus was proportional to the intensity of the reaction'' という文を挙げる。彼の与える例では、それが実際に意味しているのは、子どもたちは truck より bunny のほうを長く見る、ということにすぎない。最初の文は、必ずしも偽ではない。だが、それは場面を差し出さない。二つ目の文は、読者に見えるものを与える。

\subsection{Question \& Answer}

\paragraph{Question.} 書き手たちはしばしば知的で誠実なのに、なぜこれほど悪文が多いのか。

\paragraph{Answer.} ある分野における知性は、読者がどこにいるかを見定める技能を保証しないからである。書き手の知識が、読者についての書き手のモデルを追い越してしまうのだ。Pinker の中心的な一手は、虚栄や悪意についての物語を、視点取得の失敗についての物語へ置き換えることにある。

\begin{proposition}
説明的散文における支配的な失敗の型は、アイデアの欠如ではなく、読者が何をどこまで無理なく推論できるかを見積もり損ねることである。
\end{proposition}

だからこそ Pinker は、the curse of knowledge を egocentrism や theory of mind の欠如と結びつける。したがってこの講義は、のちに examples、構造、リズム、style をめぐって発せられるあらゆる問いを支配する、ひとつの認知的な誤りから始まるのである。

\section{読者への較正を、書くことの実践にする}

診断が述べられると、講義はただちに remedy を求める。もし悪文が読者モデルの失敗なのだとすれば、草稿のあいだに私たちは何をすべきなのか。Pinker の第一の答えは慎ましい。読者を想像してみること。共感を養うこと。だが彼は、ほとんどすぐに、それだけでは足りないと付け加える。the curse of knowledge が厄介なのは、まさにそれを内部からは感じられないからである。説明するには自明すぎると思われることが、実際にはまったく自明でないことがしばしばある。

図式的には、推敲の過程は
\begin{align}
\text{草稿} &\to \text{想像された読者}, \\
\text{草稿} + \text{実際の読者の反応} &\to \text{revision}.
\end{align}
のように見える。第一の矢印は必要である。第二の矢印が、それを救う。

Pinker がここでとりわけ優れているのは、読者についての粗雑な像を拒むところにある。彼が草稿を母親に見せたのは、彼女が鈍かったり、洗練を欠いていたからではない。むしろ逆だと彼は言う。彼女は知的で、よく読み、真面目だった。重要だったのは、彼女が彼の分野の specialist ではなかったことだ。適切な test reader とは、人口の中から無作為に選ばれた誰かでもなければ、同業者でもない。きちんと述べられれば議論を追えるだけの知性を持ちながら、書き手の局所的 shorthand は共有していない人である。

講義は、具体的な読者の系列を通っていく。
\begin{itemize}
\item 書き手自身による共感の試み
\item 信頼できる、知的な non-specialist
\item 商業出版社の editor
\item 隣接分野の academics
\item さらには、書き手のごく小さな研究圏のすぐ外に立つ学生や同僚
\end{itemize}

この最後の点は重要であり、しかも見失われやすい。Pinker は、academy の内部でさえ、ときには同じ department の中でさえ、人々は互いに理解不能になってしまうことがある、と指摘する。lab の中の五、六人が何か月も同じ jargon を吸い続けていれば、そうなる。いったん prose がその microclimate を出れば、たちまち崩れてしまう。

したがって、この講義の最初の実践的な rule は、単に ``読者を知れ'' ではない。もっと厳しい。
\begin{quote}
読者の頭の中へ入ろうとしなさい。だが、その努力だけを信用してはならない。賢く、好奇心があり、しかし当面の circle の外にいる、現実の人間の前に草稿を置きなさい。
\end{quote}

これが第一の主要な修復機構である。第二の機構は、ここから自然に続く。

\section{具体性、イメージ、そして感覚的理解}

ここで講義は、この章全体にとって重要な仕方で向きを変える。読者への較正について語ったあとで、Pinker は理解そのものが何からできているのかを問う。彼の答えは、言語は、その重要性にもかかわらず、いくぶん過大評価されている、というものである。書き手として私たちは language の中で生きているし、Pinker 自身も職業として language を研究している。だが、理解とは、ただ言語的な列が一つまた一つと受け取られることではない。language は、読者にある idea を grasp させるための手段であり、その idea は、しばしば視覚的であり、運動感覚的であり、身体的であり、感情的であり、聴覚的であり、あるいはその他の感覚的なものなのである。

この主張は、次のように圧縮できる。
\begin{equation}
\text{language} \to \text{mental model},
\qquad
\text{not merely}
\qquad
\text{language} \to \text{more language}.
\end{equation}

これが Pinker の第二の主要 rule である。第一の rule が読者の位置を定めることだとすれば、第二の rule は、読者に組み立てるべきものを与えることだ。``stimulus'' と言うな、もし意味しているのが ``bunny'' なのなら。level、perspective、framework、paradigm、concept の背後に隠れるな、もし読者に必要なのがひとつの scene なのだとしたら。

\paragraph{A worked example.}
この講義は、抽象から可視性へのほとんど完璧な小さな導出を与えてくれる。
\begin{enumerate}
\item ``the level of the stimulus was proportional to the intensity of the reaction'' という文から始める。
\item その文が、関係は記録しているが、知覚可能な内容のほとんどすべてを抑圧していることに気づく。
\item 読者は何を思い描くべきなのかを問う。
\item 具体的な対象と行為を復元する。子ども、bunny、truck、長く見ること。
\item 主張は prose として見えるようになり、したがって考えられるようになる。
\end{enumerate}

図式的には、
\begin{align}
\text{``The level of the stimulus was proportional to the intensity of the reaction.''}
&\not\Rightarrow \text{読者にとって安定した scene}, \\
\text{``Kids look longer at a bunny than at a truck.''}
&\Rightarrow \text{検分可能な mental picture}.
\end{align}
と書ける。

講義は、この一例で止まらない。visual metaphor へと広がっていく。昔の prose が私たちにいっそう vivid に感じられるのは、昔の書き手たちには、今すぐ使える abstractions が少なく、したがって共有された images により強く依拠せざるをえなかったからだ、と Pinker は示唆する。彼は、``the spirit of the hawk kneaded into our flesh'' という印象的な句を挙げ、それが、後代の technical term である ``aggression'' や ``antisocial behavior'' がいま担っている仕事を、かつては果たしていたような表現なのだと言う。私たちは古い style を丸ごと模倣する必要はない。だが、その mechanism は理解しておくべきだ。書き手は、想像できるものを差し出すことによって、読者に届くのである。

\section{なぜ書くことは話すことより難しいのか}

ここで講義は、新たな障害を持ち出す。speech は子どもにとってあれほど自然に生じるのに、なぜ writing はなお遅く、骨が折れ、不自然なのか。Pinker は、レベルを変えることによって答える。ここまでは advice を与えてきた。ここでは、その advice が必要になる環境そのものを説明する。

\subsection{Question \& Answer}

\paragraph{Question.} speech に比べて、なぜ writing は不自然に感じられるのか。

\paragraph{Answer.} conversation は共有された situation の内部で始まるのに対し、writing は、ページ上の marks だけから自分自身の situation を構築しなければならないからである。

講義の対比は、やはり図式的に、次のように書ける。
\begin{align}
\text{conversation} &= \text{common ground} + \text{deixis} + \text{feedback}, \\
\text{writing} &= \text{ページ上のテクスト}.
\end{align}

この小さな二つの式は、transcript にあるいくつかの異なる点を凝縮している。
\begin{enumerate}
\item conversation は、無から始まるのではない。参加者たちは、なぜそこにいて、どんな occasion が自分たちを引き合わせたのかを知っている。
\item context が共有されているので、``this''、``that''、``she''、``the thing''、``what I was talking about'' のような語も、しばしば完全に明瞭である。
\item conversation には即座の correction がある。眉のひそめ、困惑した目つき、説明を求める言葉、注意の漂い。
\item live public speaking でさえ、この支えのいくらかをなお保っている。話し手が聴衆を見られるからである。
\item writing は、その支えを剥ぎ取る。読者は別の場所に住み、後に読み、prose が書かれた局所的 setting について何も知らないかもしれない。
\end{enumerate}

講義の rhythm は、この拍に依存している。これがなければ、examples や明示的 framing は、任意のもの、あるいは装飾のように聞こえてしまうかもしれない。これがあることで、それらは構造的な必需品になる。writing が speech より難しいのは、conversation では situation、gesture、feedback が担っている burden を、page が一人で担わなければならないからである。

\section{一般化、例、そして意味のかたち}

ここで講義は、大きな対比から、私たちが一文ごとに用いうる方法へと焦点を狭める。聞き手は、一つの principle を、あえて鋭い形で提案する。例を欠く generalization も、generalization を欠く例も、どちらも役に立たない、と。Pinker はこの verdict を少し和らげるが、その point の構造自体は受け入れる。case を欠いた generalization は、読者に、何が主張されているのかを本当には知らぬまま頷かせてしまうことが多い。generalization を欠いた example は、そもそもなぜその例が導入されたのかと読者に問い返させてしまう。

この均衡は、
\begin{equation}
\text{examples} + \text{generalizations} \to \text{理解}.
\end{equation}
と要約できる。聞き手はそのあと、同じ均衡を context と compression のあいだのものとして言い換える。これは Pinker 自身の厳密な terminology ではないが、このやりとりを要約する編集上のまとめとしては有用である。
\begin{equation}
\text{context} \leftrightarrow \text{compression}.
\end{equation}

論理は次のような段階で展開される。
\begin{enumerate}
\item generalization は圧縮する。それは particulars を横断して一掃する。
\item 細部を抑圧するがゆえに、それは確定した referent を呼び起こし損ねることがある。
\item example は、具体的な ballpark を回復する。
\item だが example だけでは、まだ何が学ばれたのかは語られない。
\item generalization が戻ってきて、その example がなぜ重要なのかを告げる。
\end{enumerate}

ここで講義は、ふつうの writing lesson より微妙なことを行っている。exposition から semantics へと向きを変えるのである。Pinker は、いくつかの familiar な expressions を挙げるが、それらの意味は、それを構成する部分の意味をただ足し合わせれば回復できるものではない。bathroom は、bath を含んでいる必要はない。going to the bathroom は、文字どおりには浴槽のある room へ行くことを意味しない。breakfast は、fast を破ることとして逐一経験される必要はない。Christmas も、それが utter されるたびに、あらためて compositional theology の新鮮な一片として処理されているわけではない。

そして講義は、さらに広がる。olive oil は olives から作られた oil だが、baby oil は babies のための oil である。同じ surface form が、異なる semantic relation を支えている。Richard Lederer の、より playful な examples は、この lesson をさらに押し進める。Singer は、``-er'' が verb を typical agent の noun に変えるという生きた rule を、透けて見せている。Finger はそうではない。language は、元の logic が生きた意識から消えたあとも、古い formations を保存し続けるのである。

\begin{remark}
意味は、しばしば機械的に compositional なのではなく、conventional であり、historical である。Pinker の point は、rule がまったく存在しないということではない。読者は、surface form だけから intended meaning を回復できるはずだ、と期待してはならない、ということだ。
\end{remark}

講義は短く adultery と adulterate にも触れ、semantic drift の一例として示す。この moment を notes に保存する prudent な仕方は、その example に講義以上の重さを担わせないことである。重要なのは、familiar な words が元の関係から大きく離れていくことがありうるという点、そして書き手は、formal decomposition だけで理解が得られると決して仮定してはならない、という点である。

\section{快楽としての style: rhythm、sound、brevity、humor}

講義は、意味が words によって自動的には運ばれないことを示したあとで、ようやく style へ向かう。Pinker は順序に慎重である。彼は快楽を clarity の代用品として提案するのではない。clarity が確保されたあとにくる、次の層として提案するのである。

講義が Shakespeare と Strunk へ至る前に、子どもの説明の freshness に立ち止まる。Pinker は、``smoke is fire vapor'' のような example を喜ぶ。なぜなら、それは inherited cliches ではなく、見える関係へ向かってなお手を伸ばしている mind を示しているからである。子どもたちは、後の書き手たちが refuge として用いる重い abstractions の inventory を、まだ持っていない。その意味で、彼らは、後代の technical vocabulary に頼れなかったため共有可能な imagery へ手を伸ばさざるをえなかった、昔の書き手たちに似ている。lesson は、子どものように書けということではない。freshness と observability は、深く結びついている、ということである。

\subsection{Question \& Answer}

\paragraph{Question.} prose を、単に理解可能なだけでなく、快いものにしているものは何か。

\paragraph{Answer.} Pinker の答えでは、prose には audible な生命がある。うまく動くこともあれば、まずく動くこともある。rhythm を運ぶこともあれば、それを壊すこともある。耳に障ることもあれば、滑るように流れることもある。書き手はこれを、read aloud し、friction に耳を澄まし、その line がきれいに spoken されるまで revise することによって、直接試すことができる。

局所的な update rule は単純である。
\begin{equation}
\text{read aloud} \to \text{friction を聞き取る} \to \text{cadence を revise する}.
\end{equation}

いくつかの stylistic variables が、曖昧な称賛ではなく、操作可能な quantities としてここに入ってくる。
\begin{itemize}
\item \textbf{visual imagery}: その sentence は scene を呼び起こせるか。
\item \textbf{metrical structure}: prose は、つまずくのではなく、使える beats で進んでいるか。
\item \textbf{sibilance}: 連続する hiss 音が多すぎないか。
\item \textbf{alliteration}: 無理に感じられるほどではなく、しかし軽い快楽を与えるだけの sound の反復があるか。
\item \textbf{brevity}: 同じ意味を、より少ない drag で届けられるか。
\end{itemize}

Pinker は humor を、この部分に直接結びつける。humor は freshness と timing に依存する。joke は引き延ばしすぎれば funny でなくなる。punch line は、過剰に予告されれば、その preparation 自体の重みによって潰れてしまう。したがって \emph{Hamlet} の ``brevity is the soul of wit'' という line は、二重に働く。wit についての真理を述べると同時に、その economy によってその真理を exemplify しているのである。

Strunk の ``omit needless words'' も、同じ仕方で扱われる。それは handbook から引用された slogan にすぎないのではない。それ自体の example でもあり、Pinker は自分の writing life から practical な test を付け加える。新聞や magazine が hard word limit を課すと、その強制された compression は、単に prose を短くするだけでなく、しばしばそれを改善するのである。

図式的には、
\begin{equation}
\text{needless words} \to \text{余分な processing} \to \text{弱くなる force}.
\end{equation}
となる。

ここでの講義の argument は、認知的でもあり、美的でもある。余分な一語一語は、すべて読者にとって余分な work である。だがそれ以上に、compression はしばしば rhythm を鋭くし、書き手を woolly な idiom から concrete phrasing へ押しやる。この説明において、style とは discipline された pleasure なのである。

\section{古い prose、academic prose、そして LLM prose}

ここで講義は、craft から institutions、history、technology へと広がっていく。Pinker はまず、なぜ自分が bad academic prose にこれほど強い関心を持つのかを説明する。その complaint は、単なる stylistic fastidiousness ではない。bad prose は intelligence を浪費する。公共的価値を無駄にする。reviewers、colleagues、students に、本来なら source の段階で明らかにされるべき sentence を decode するための時間を費やさせる。claim が現れるまで五回も六回も読まなければならない paragraph は、単に醜いだけではない。非効率であり、危険でもある。

その institutional complaint は、ついで historical な complaint へと開いていく。なぜ older prose は、現代の読者に、これほど vivid に、あるいは豊潤にさえ感じられるのか。Pinker は、いくつかの理由を順に挙げる。昔の書き手たちは、強い literary models によって教育されていた。prose は、世界に対して自分を示す medium として、より公然と cultivate されていた。そして、この講義の内的 logic にとって最も重要なのは、昔の書き手たちには ready-made な abstractions や cliches が少なく、それゆえ ordinary readers に image を結ばせられる forms の中へ、新しい ideas を cast せざるをえなかった、ということである。

その historical argument は、さらに cultural change へと折れ曲がる。Pinker は、informalization の長い process に名前を与える。hierarchy が減り、ceremony が減り、conversational intimacy が増し、elevated あるいは self-consciously polished に聞こえるものへの suspicion が増す。この shift は、強く crafted な prose を書ける modern writers でさえ、formal すぎたり、仕上がりすぎて聞こえたりすることをためらう理由の説明になる。

こうしたことを経てはじめて、講義は AI へ戻る。この順序が重要である。LLM prose は、異質な object として導入されるのではない。すでに存在していた problem に対する、最新の pressure として導入される。Pinker はまず、その strengths を認める。limited な意味では、それはしばしば clear である。sentences は orderly で、syntax は plain で、paragraphs は recognizable な openings と closings を持ちがちだからである。だがその weakness も同じく際立っている。それは generic であり、prosaic であり、aggregate style の pastiche としてしばしば一目で分かる。

この対比は、図式的には次のように述べられる。
\begin{align}
\text{LLM prose} &= \text{clear な ordering} + \text{plain な syntax} + \text{generic な pastiche}, \\
\text{強い prose} &= \text{clarity} + \text{concreteness} + \text{freshness}.
\end{align}

したがって講義は、ひとつの重要な non-equivalence をもって閉じる。
\begin{equation}
\text{clear な prose} \neq \text{fresh な prose}.
\end{equation}

Pinker は短く、なぜ LLM prose が academics や lawyers や bureaucrats の prose よりしばしば読みやすいのかを推測する。それはおそらく training と feedback によって hammer で形を整えられてきたのだろう。あるいは、より speculative に言えば、無数の examples の composite が、もっともひどい局所的 distortions のいくつかを取り除いているのかもしれない。だが彼は、その speculation を、講義が支えられる以上の強い claim へと漂わせはしない。人間の mind は単なる大規模言語モデルなのだ、とは結論しない。むしろ、より狭い lesson を引き出す。intelligence についての serious な account は、いまや older rule-centered views がそうした以上に、大規模入力からの大域的な pattern extraction に大きな重みを与えねばならない、ということである。

\section{まとめ}

この講義は連鎖として展開する。そして、そのように読むのが最もよい。まず、bad prose はおもに戦略的な obscurity なのだという、気分のよい物語を退ける。代わりに the curse of knowledge を据え、ただちに修復へ向かう。現実の readers、現実の editors、現実の tests である。そこからさらに議論を深め、理解とは何かを問う。そして、読者に必要なのは words だけではなく mental models なのだ、と答える。この答えは、speech と writing の非対称へつながり、そこから example と generalization の均衡へつながる。意味は自動的なものでも、完全に compositional なものでもないと示したあとで、講義はようやく style を operational なものとして扱うことができるようになる。rhythm、sound、compression、wit、pleasure である。最後に講義は history、institutions、AI へと広がるが、最初の point を見失わない。最初のページから最後のページまで、書き手の task は変わらない。自分たちの意味することを、他の mind が実際に入り込める form で述べることなのである。